\documentclass[11pt]{article} 
\usepackage[margin=1in]{geometry}
\usepackage{amsfonts,latexsym,amsthm,amssymb,amsmath,amscd,euscript,esint, dsfont,mathtools}	
\usepackage{graphicx}		
\usepackage{hyperref}
\usepackage{cases}			
\usepackage{textcomp, gensymb}
\usepackage{xcolor}
\usepackage{tabularx}

% diagrams
\usepackage{quiver}

\makeatletter
\def\namedlabel#1#2{\begingroup
	\def\@currentlabel{#2}
	\label{#1}\endgroup
}
\makeatother

\setlength{\parindent}{0pt} 	% first line in paragraph will not be indented

\usepackage[parfill]{parskip}
\usepackage{lastpage}
\usepackage{fancyhdr}
\newcommand{\ind}{{\mathds{1}}}

% math fonts
\renewcommand{\AA}{\mathbb{A}}
\newcommand{\BB}{\mathbb{B}}
\newcommand{\CC}{\mathbb{C}}
\newcommand{\DD}{\mathbb{D}}
\newcommand{\EE}{\mathbb{E}}
\newcommand{\FF}{\mathbb{F}}
\newcommand{\GG}{\mathbb{G}}
\newcommand{\HH}{\mathbb{H}}
\newcommand{\II}{\mathbb{I}}
\newcommand{\JJ}{\mathbb{J}}
\newcommand{\KK}{\mathbb{K}}
\newcommand{\LL}{\mathbb{L}}
\newcommand{\MM}{\mathbb{M}}
\newcommand{\NN}{\mathbb{N}}
\newcommand{\OO}{\mathbb{O}}
\newcommand{\PP}{\mathbb{P}}
\newcommand{\QQ}{\mathbb{Q}}
\newcommand{\RR}{\mathbb{R}}
\renewcommand{\SS}{\mathbb{S}}
\newcommand{\TT}{\mathbb{T}}
\newcommand{\UU}{\mathbb{U}}
\newcommand{\VV}{\mathbb{V}}
\newcommand{\WW}{\mathbb{W}}
\newcommand{\XX}{\mathbb{X}}
\newcommand{\YY}{\mathbb{Y}}
\newcommand{\ZZ}{\mathbb{Z}}

\newcommand{\cA}{\mathcal{A}}
\newcommand{\cB}{\mathcal{B}}
\newcommand{\cC}{\mathcal{C}}
\newcommand{\cD}{\mathcal{D}}
\newcommand{\cE}{\mathcal{E}}
\newcommand{\cG}{\mathcal{G}}
\newcommand{\cF}{\mathcal{F}}
\newcommand{\cH}{\mathcal{H}}
\newcommand{\cI}{\mathcal{I}}
\newcommand{\cJ}{\mathcal{J}}
\newcommand{\cK}{\mathcal{K}}
\newcommand{\cL}{\mathcal{L}}
\newcommand{\cM}{\mathcal{M}}
\newcommand{\cN}{\mathcal{N}}
\newcommand{\cO}{\mathcal{O}}
\newcommand{\cP}{\mathcal{P}}
\newcommand{\cQ}{\mathcal{Q}}
\newcommand{\cR}{\mathcal{R}}
\newcommand{\cS}{\mathcal{S}}
\newcommand{\cT}{\mathcal{T}}
\newcommand{\cU}{\mathcal{U}}
\newcommand{\cV}{\mathcal{V}}
\newcommand{\cW}{\mathcal{W}}
\newcommand{\cX}{\mathcal{X}}
\newcommand{\cY}{\mathcal{Y}}
\newcommand{\cZ}{\mathcal{Z}}

\newcommand{\ka}{\mathfrak{a}}
\newcommand{\kb}{\mathfrak{b}}
\newcommand{\kc}{\mathfrak{c}}
\newcommand{\kd}{\mathfrak{d}}
\newcommand{\ke}{\mathfrak{e}}
\newcommand{\kg}{\mathfrak{g}}
\newcommand{\kf}{\mathfrak{f}}
\newcommand{\kh}{\mathfrak{h}}
\newcommand{\ki}{\mathfrak{i}}
\newcommand{\kj}{\mathfrak{j}}
\newcommand{\kk}{\mathfrak{k}}
\newcommand{\kl}{\mathfrak{l}}
\newcommand{\km}{\mathfrak{m}}
\newcommand{\kn}{\mathfrak{n}}
\newcommand{\ko}{\mathfrak{o}}
\newcommand{\kp}{\mathfrak{p}}
\newcommand{\kq}{\mathfrak{q}}
\newcommand{\kr}{\mathfrak{r}}
\newcommand{\ks}{\mathfrak{s}}
\newcommand{\kt}{\mathfrak{t}}
\newcommand{\ku}{\mathfrak{u}}
\newcommand{\kv}{\mathfrak{v}}
\newcommand{\kw}{\mathfrak{w}}
\newcommand{\kx}{\mathfrak{x}}
\newcommand{\ky}{\mathfrak{y}}
\newcommand{\kz}{\mathfrak{z}}

\newcommand{\kA}{\mathfrak{A}}
\newcommand{\kB}{\mathfrak{B}}
\newcommand{\kC}{\mathfrak{C}}
\newcommand{\kD}{\mathfrak{D}}
\newcommand{\kE}{\mathfrak{E}}
\newcommand{\kG}{\mathfrak{G}}
\newcommand{\kF}{\mathfrak{F}}
\newcommand{\kH}{\mathfrak{H}}
\newcommand{\kI}{\mathfrak{I}}
\newcommand{\kJ}{\mathfrak{J}}
\newcommand{\kK}{\mathfrak{K}}
\newcommand{\kL}{\mathfrak{L}}
\newcommand{\kM}{\mathfrak{M}}
\newcommand{\kN}{\mathfrak{N}}
\newcommand{\kO}{\mathfrak{O}}
\newcommand{\kP}{\mathfrak{P}}
\newcommand{\kQ}{\mathfrak{Q}}
\newcommand{\kR}{\mathfrak{R}}
\newcommand{\kS}{\mathfrak{S}}
\newcommand{\kT}{\mathfrak{T}}
\newcommand{\kU}{\mathfrak{U}}
\newcommand{\kV}{\mathfrak{V}}
\newcommand{\kW}{\mathfrak{W}}
\newcommand{\kX}{\mathfrak{X}}
\newcommand{\kY}{\mathfrak{Y}}
\newcommand{\kZ}{\mathfrak{Z}}

%some math symbol commands redefined:
\DeclareMathOperator{\C}{\mathcal{C}}
\DeclareMathOperator{\power}{\mathcal{P}}
\DeclareMathOperator{\vspan}{\text{span}}
\DeclareMathOperator{\vnull}{\text{null}}
\DeclareMathOperator{\range}{\text{range}}
\DeclareMathOperator{\re}{\text{Re}}
\DeclareMathOperator{\im}{\text{Im}}
\DeclareMathOperator{\erf}{\text{erf}}
\newcommand{\pdv}[2]{\frac{\partial #1}{\partial #2}}
\newcommand{\pdvv}[2]{\frac{\partial^2 #1}{\partial #2}}
\DeclareMathOperator{\tr}{\operatorname{Tr}}
\newcommand{\Deck}{\operatorname{Deck}}
\newcommand{\Conj}{\mathrm{Conj}}
\newcommand{\Sing}{\mathrm{Sing}}
\DeclareMathOperator{\Spec}{\operatorname{Spec}}
\newcommand{\Proj}{\operatorname{Proj}}
\newcommand{\Res}{\operatorname{Res}}
\newcommand{\PROJ}{\mathit{Proj}}

% category names
\newcommand{\Set}{\mathsf{Set}}
\newcommand{\Grp}{\mathsf{Grp}}
\newcommand{\Ring}{\mathsf{Ring}}
\newcommand{\Top}{\mathsf{Top}}
\newcommand{\Sch}{\mathsf{Sch}}

\newcommand{\ob}{\operatorname{Ob}}
\newcommand{\Id}{\operatorname{Id}}
\newcommand{\Hom}{\operatorname{Hom}}
\newcommand{\colim}{\operatorname{colim}}
\newcommand{\Ann}{\operatorname{Ann}}
\newcommand{\Supp}{\operatorname{Supp}}
\newcommand{\Ass}{\operatorname{Ass}}
\newcommand{\pre}{\text{pre}}
\newcommand{\adj}{\text{adj}}
\newcommand{\Ker}{\operatorname{Ker}}
\newcommand{\Coker}{\operatorname{Coker}}
\newcommand{\Nil}{\operatorname{Nil}}
\newcommand{\Char}{\operatorname{char}}
\newcommand{\Adj}{\operatorname{Adj}}
\newcommand{\Bl}{\mathrm{Bl}}
\newcommand{\codim}{\mathrm{codim}}
\newcommand{\val}{\operatorname{val}}
\newcommand{\Div}{\operatorname{div}}
\newcommand{\Pic}{\operatorname{Pic}}
\newcommand{\Weil}{\operatorname{Weil}}
\newcommand{\Cl}{\operatorname{Cl}}
\newcommand{\Sym}{\operatorname{Sym}}

\newcommand{\GL}{\operatorname{GL}}
\newcommand{\PGL}{\operatorname{PGL}}
\newcommand{\SL}{\operatorname{SL}}
\newcommand{\PSL}{\operatorname{PSL}}
\newcommand{\Rk}{\operatorname{Rk}}

\newtheorem{lemma}{Lemma}
\newtheorem{claim}{Claim}

% category theory symbols
\DeclareMathOperator{\Mor}{\operatorname{Mor}}

\newenvironment{solution}{\begin{trivlist}\item[]{\textcolor{blue}{Solution:}}}{\qed \end{trivlist}}


\usepackage[shortlabels]{enumitem}
\title{MATH 2060 Homework 3}
\author{Zhengyu Zou}
\date{\today}

\begin{document}
\maketitle

\section*{FOAG 17.3.F.}

\begin{solution}
    We will construct the map $\psi: X \rightarrow \PROJ \Sym^{\bullet} (\pi_* \cL)$ on the affines of $X$, and show that these maps agree with each other on the overlaps. Let $\Spec A \subset Y$ be an affine on which $\pi_* \cL$ is trivializable. The preimage of $\Spec A$ in $\PROJ \Sym^{\bullet} (\pi_* \cL)$ is given by $\Proj_A \Sym^{\bullet} (\pi_* \cL)(\Spec A)$. Since $\pi_* \cL$ is finite type, $(\pi_* \cL)(\Spec A)$ is a finitely-generated $A$-module, so $(\pi_* \cL)(\Spec A) \cong A^{\oplus n} / I'$ for some submodule $I' \subset A^{\oplus n}$. In particular, since $\pi_* \cL$ is trivializable on $\Spec A$, we have an isomorphism $A^{\oplus n} / I \rightarrow A$ sending the generator of each direct summand to a regular functions $f_i \in A$. It follows that $\Sym^{\bullet} (\pi_* \cL)(\Spec A) \cong A[x_1, \ldots, x_n] / I$, where $I$ is the submodule in $A[x_1, \ldots, x_n]$ generated by $I'$. Now, given any $\Spec B \subset X$ such that the image of $\Spec B$ under $\pi$ is contained in $A$, the pullback sheaf $\pi^*\pi_* \cL$ on $\Spec B$ is given by $B \otimes_A A[x_1, \ldots, x_n] / I \cong B[y_1, \ldots, y_n] / \pi^{\sharp}(I)$, where $\pi^{\sharp}(I)$ is the submodule generated by the image of $I$ under the map of structure ring $\pi^{\sharp}: A \rightarrow B$. Then, since $\cL$ is relatively base-point-free, we have a surjection $\pi^* \pi_* \cL \rightarrow \cL$, so $\cL(B) \cong B[x_1, \ldots, x_n] / J$ for some $J \supset \pi^{\sharp}(I)$. Here, the sections $y_i$ are the pullbacks of the sections $x_i$ on $\pi_* \cL$, and one can check that $\pi_*(y_i) = x_i$, so the $y_i$'s are base-point free. It follows that we have a map $\Spec B \rightarrow \PP_A^{n-1}$ given by the line bundle $\cL$ and the sections $y_1, \ldots, y_n$. In particular, since the sections $y_1, \ldots, y_n$ satisfy the relations in $J$, the pushforwards of the sections satisfy relations in $I$, so the image of the map $\Spec B \rightarrow \PP_A^{n-1}$ lies in $\Proj_A A[x_1, \ldots, x_n] / I$. This gives us the desired map on $X$. 
    % https://q.uiver.app/#q=WzAsNSxbMiwxLCJZIl0sWzEsMSwiXFxTcGVjIEEiXSxbMiwwLCJcXFBST0ogXFxTeW1ee1xcYnVsbGV0fSBcXHBpXyogXFxjTCJdLFsxLDAsIlxcUHJval9BIFxcZnJhY3tBW3hfMSwgXFxsZG90cywgeF9uXX17SX0iXSxbMCwxLCJcXFNwZWMgQiBcXHN1YnNldCBYIl0sWzEsMCwiIiwwLHsic3R5bGUiOnsidGFpbCI6eyJuYW1lIjoiaG9vayIsInNpZGUiOiJ0b3AifX19XSxbMywxXSxbMywyLCIiLDIseyJzdHlsZSI6eyJ0YWlsIjp7Im5hbWUiOiJob29rIiwic2lkZSI6InRvcCJ9fX1dLFsyLDBdLFs0LDFdLFs0LDMsIlxcZXhpc3QiXV0=
    \[\begin{tikzcd}
        & {\Proj_A \frac{A[x_1, \ldots, x_n]}{I}} & {\PROJ \Sym^{\bullet} (\pi_* \cL)} \\
        {\Spec B \subset X} & {\Spec A} & Y
        \arrow[hook, from=1-2, to=1-3]
        \arrow[from=1-2, to=2-2]
        \arrow[from=1-3, to=2-3]
        \arrow["\exists", from=2-1, to=1-2]
        \arrow[from=2-1, to=2-2]
        \arrow[hook, from=2-2, to=2-3]
    \end{tikzcd}\]
\end{solution}

\newpage

\section*{FOAG 18.1.A}

\begin{solution}
    Let $\cF$ be an arbitrary coherent sheaf on $X$. Since $Y$ is locally Noetherian, it suffices to show that $\pi_* \cF$ is finite type. Since $\pi: X \rightarrow Y$ is projective, there exists a quasicoherent sheaf of graded $\cO_Y$-algebras $\cI_{\bullet}$ such that $X \cong \PROJ \cI_{\bullet}$ and the map $\pi$ factors through the isomorphism as 
    % https://q.uiver.app/#q=WzAsMyxbMCwwLCJYIl0sWzEsMSwiWSJdLFsyLDAsIlxcUFJPSiBcXGNJX3tcXGJ1bGxldH0iXSxbMCwxLCJcXHBpIiwyXSxbMCwyLCJcXHNpbWVxIl0sWzIsMV1d
    \[\begin{tikzcd}
        X && {\PROJ \cI_{\bullet}} \\
        & Y
        \arrow["\simeq", from=1-1, to=1-3]
        \arrow["\pi"', from=1-1, to=2-2]
        \arrow[from=1-3, to=2-2]
    \end{tikzcd}\]
    Therefore, we may identify $X$ with $\PROJ \cI_{\bullet}$ and treat $\cF$ to be a coherent sheaf on $\PROJ \cI_{\bullet}$.

    We would like to show that on each affine $\Spec A \subset Y$, the sheaf $\pi_* \cF|_{\Spec A}$ is a finitely-generated $A$-module. Recall by the construction of global Proj that the preimage of the affine $\Spec A$ is given by $\Proj_A \cI_{\bullet}(\Spec A)$. Since $\cI_{\bullet}$ is quasicoherent and finitely generated in degree-1, we may identify $\cI_{\bullet}(\Spec A)$ with $A[x_1, \ldots, x_n] / I$. We then have the following diagram:
    % https://q.uiver.app/#q=WzAsNCxbMCwwLCJcXFByb2pfQSBcXGZyYWN7YVt4XzEsIFxcbGRvdHMsIHhfbl19e0l9Il0sWzEsMCwiXFxQUk9KIFxcY0lfe1xcYnVsbGV0fSJdLFswLDEsIlxcU3BlYyBBIl0sWzEsMSwiWSJdLFsyLDMsIiIsMCx7InN0eWxlIjp7InRhaWwiOnsibmFtZSI6Imhvb2siLCJzaWRlIjoidG9wIn19fV0sWzEsM10sWzAsMl0sWzAsMV1d
    \[\begin{tikzcd}
        {\Proj_A \frac{A[x_1, \ldots, x_n]}{I}} & {\PROJ \cI_{\bullet}} \\
        {\Spec A} & Y
        \arrow[from=1-1, to=1-2]
        \arrow[from=1-1, to=2-1]
        \arrow[from=1-2, to=2-2]
        \arrow[hook, from=2-1, to=2-2]
    \end{tikzcd}\]
    Since $\cF$ is coherent, $\cF|_{\pi^{-1}(\Spec A)}$ is a finite type quasicoherent sheaf over $\Proj_A \frac{A[x_1, \ldots, x_n]}{I}$. Then, Theorem 16.1.1 implies there exists a surjection of modules 
    % https://q.uiver.app/#q=WzAsMixbMCwwLCJcXGJpZ29wbHVzX3trPTF9Xm4gXFxjT197XFxQcm9qX0EgXFxmcmFje0FbeF8xLCBcXGxkb3RzLCB4X25dfXtJfX0oLW0pICJdLFsxLDAsIlxcY0YoXFxwaV57LTF9KFxcU3BlYyBBKSkiXSxbMCwxLCIiLDAseyJzdHlsZSI6eyJoZWFkIjp7Im5hbWUiOiJlcGkifX19XV0=
    \[\begin{tikzcd}
        {\bigoplus_{k=1}^n \cO_{\Proj_A \frac{A[x_1, \ldots, x_n]}{I}}(-m) } & {\cF|_{\pi^{-1}(\Spec A)}}
        \arrow[two heads, from=1-1, to=1-2]
    \end{tikzcd}\]
    Here, since $\cO(-m)$ is a finitely-generated $A$-module, so is a finite direct sum of $\cO(-m)$. This implies $\cF|_{\pi^{-1}(\Spec A)}$ is a finitely-generated $A$-module, so $\pi_* \cF$ is finite-type. Finally, since $Y$ is locally noetherian, we are done.
\end{solution}

\newpage

\section*{FOAG 18.1.B.}

\begin{solution}
    We will use the result from 12.3.G that if $X$ and $Y$ are pure-dimensional closed subvarieties of $\PP^n$ of codimensions $d$ and $e$ respectively, and $d + e \leq n$, then $X$ and $Y$ intersect. 

    Given any $q \in Y$, we would like to find an open $U \ni q$ such that for all $p \in U$, $\dim \pi^{-1}(p) \leq \dim \pi^{-1}(q)$. Let $\Spec A \subset Y$ be a distinguished affine open that contains $q$. The morphism $\pi: \pi^{-1}(\Spec A) \rightarrow \Spec A$ is again projective, so we may assume $X = \pi^{-1}(\Spec A)$ and $Y = \Spec A$. We can then factor the morphism as 
    % https://q.uiver.app/#q=WzAsMyxbMCwwLCJYIl0sWzEsMCwiXFxQUF9BXm4iXSxbMSwxLCJcXFNwZWMgQSJdLFswLDIsIlxccGkiLDJdLFswLDEsIiIsMCx7InN0eWxlIjp7InRhaWwiOnsibmFtZSI6Imhvb2siLCJzaWRlIjoidG9wIn19fV0sWzEsMl1d
    \[\begin{tikzcd}
        X & {\PP_A^n} \\
        & {\Spec A}
        \arrow[hook, from=1-1, to=1-2]
        \arrow["\pi"', from=1-1, to=2-2]
        \arrow[from=1-2, to=2-2]
    \end{tikzcd}\]
    where the morphism $X \hookrightarrow \PP_A^n$ is a closed embedding. If $q$ is the generic point of $\Spec A$, then for all $p \in \Spec A$ we must have $\pi^{-1}(p) \subset \pi^{-1}(q) = X$, so $\dim \pi^{-1}(p) \leq \dim \pi^{-1}(q)$, and we may just take $U = \Spec A$. Otherwise, $\pi^{-1}(q)$ is a proper subset of $\PP_A^n$. Write $r = \dim \pi^{-1}(q)$. From 12.3.D, there exists hypersurfaces $H_1, \ldots, H_{r+1}$ in $\PP_{\kappa(q)}^n$ such that their intersection doesn't contain points in $\pi^{-1}(q)$. In particular, 12.3.G implies the irreducible components of $\bigcap_i H_i$ must all have codimension $r$. Lifting the defining equations to $A$, we obtain hypersurfaces $H'_1, \ldots, H'_{r+1}$ in $\PP_A^n$ whose intersection doesn't intersect $\pi^{-1}(q)$. Consider the closed set $X \cap (\bigcap_i H'_i)$. By the fundamental theorem of elimination theory, $H' = \pi(X \cap (\bigcap_i H'_i))$ is closed in $\Spec A$. Let $U = \Spec A - H'$. Then, we have a new diagram 
    % https://q.uiver.app/#q=WzAsMyxbMCwwLCJcXHBpXnstMX0oVSkiXSxbMSwwLCJcXFBQX0FebiBcXHNldG1pbnVzIEgiXSxbMSwxLCJVIl0sWzAsMiwiXFxwaSIsMl0sWzAsMSwiIiwwLHsic3R5bGUiOnsidGFpbCI6eyJuYW1lIjoiaG9vayIsInNpZGUiOiJ0b3AifX19XSxbMSwyXV0=
    \[\begin{tikzcd}
        {\pi^{-1}(U)} & {\PP_A^n \setminus \bigcap_i H'_i} \\
        & U
        \arrow[hook, from=1-1, to=1-2]
        \arrow["\pi"', from=1-1, to=2-2]
        \arrow[from=1-2, to=2-2]
    \end{tikzcd}\] 
    We claim that $U$ is the desired set. Given $p \in U$, consider the image of each hypersurface $H'_i$ in $\PP_{\kappa(p)}^n$. Call these hypersurfaces $H_{i,p}$. The intersection $\bigcap_i H_{i,p}$ contains a pure-dimensional closed subvariety of $\PP_{\kappa(p)}^n$ of codimension $r$. Then, if $\pi^{-1}(q)$ has dimension $r+1$, one of its pure-dimensional components must have codimension $< n-r$, but then by 12.3.G this component must intersect $\bigcap_i H_{i,p}$, contradicting $p \in U$. This completes the proof. 
\end{solution}

\newpage

\section*{FOAG 18.1.C.}

\begin{solution}
    Recall from Theorem 18.1.3 that $H^1(X, \cF(n)) = 0$ for $n$ large enough. Since $\cF(n) \cong \cF \otimes \cO(n)$, and that tensoring with a locally free sheaf preserves exactness, the following sequence is again exact:
    \[
        0 \longrightarrow \cF(n) \longrightarrow \cG(n) \longrightarrow \cH(n) \longrightarrow 0.
    \]
    Applying the long exact sequence of cohomology gives us the following long exact sequence 
    \[
        0 \longrightarrow H^0(X, \cF(n)) \longrightarrow H^0(X, \cG(n)) \longrightarrow H^0(X, \cH(n)) \longrightarrow H^1(X, \cF(n)) \cong 0.
    \]
    This completes the proof. 
\end{solution}

\newpage

\end{document}